\scanpage{99}
$0_{\mathrm{IV}}$\quad\uline{Différentielles et dérivations.}

\deleted{NB} Il nous a paru utile de développer les notions différentielles en utilisant la notion d'algèbre topologique formellement simple,définie dans ``Cohen''. Dans les premiers numéros qui suivent, nous n'utiliserons que la définition ``Cohen''1.1. et éventuellement une partie des sorites contenus dans 1.2.

NB Il se trouvera donc que cette rédaction contient un bon nombre de résultats de ma rédaction ``Cohen différentiel'', qu'il y aura lieu de refondre en conséquence.

1. \uline{Dérivations et extensions d'algèbres}.

\uline{Proposition} 1.1. Soient $A$ un anneau, $B$ une $A$-algèbre, $p:E\quad C$ un homomorphisme de $A$-algèbres dont le noyau $J$ est de carré nul, $u:B\quad C$ un homomorphisme de $A$-algèbres. Supposons qu'on ait un homomorphisme de $A$-algèbres $v_1:A\quad E$ qui relève $u$ i.e. tel que $pv_1=u$. Pour qu'une application $v_2:B\quad E$ ait la même propriété, il f et s qu'on ait
\[v_2-v_1=D\quad\text{i.e.}\quad v_2=v_1+D\]
\margin{$\unsure{\supset}$}
où $D:B\quad J$ est un application satisfaisant les deux conditions suivantes: $1^\circ$ $D$ est $A$-linéaire $2^\circ$ pour $f,g\quad B$, on a $D(fg)=fDg+D(f)g$
\margin{$\unsure{\supset}$}

On notera que compte tenu de $2^\circ$), $1^\circ$) équivaut aussi à \deleted{1 bis}

1 bis) $D$ s'annule sur $\operatorname{Im}(A\quad B)$.

\uline{Définition} 1.2. D'une dérivation d'un anneau dans un module, d'une $A$-dérivation d'une $A$-algèbre $B$ dans un $B$-module $J$.

\uline{Corollaire} 1.3. Etant donnés $A,B$, \deleted{$E$} $p:E\quad C,J$, $u:B\quad C$ comme dans 1.1. \deleted{\unsure{Etant ainsi} pour tout relèvement} l'ensemble des relèvements de $u$ en un homomorphisme $v$ de $A$-algèbres est vide, ou un ensemble principal homogène sous l'action du groupe $\operatorname{Dér}_A(B,J)$
\margin{98}
